\documentclass[10pt]{article}

\usepackage[margin=0.82in]{geometry}
\usepackage{amsmath,amssymb}
\usepackage{booktabs}
\usepackage{enumitem}
\usepackage{microtype}

\setlist{nosep,leftmargin=*}
\setlength{\parskip}{0.45em}
\setlength{\parindent}{0pt}

\title{\textbf{How Stable Are Regression Conclusions Across Subpopulations?}\\
\large A Brief for Preprint Review}
\author{Research summary}
\date{October 2026}

\begin{document}
\maketitle

\section*{Question and motivation}

Applied researchers commonly fit one regression to a pooled sample and interpret the resulting coefficients as if they describe every relevant subgroup. This is efficient, but it can be misleading. A pooled coefficient may conceal different subgroup relationships, and a model that performs acceptably on average may not transfer to every group.

This project asks a focused diagnostic question: \emph{when does one pooled linear regression provide a stable summary of relationships observed across predefined subpopulations?} The aim is not to introduce a new estimator or replace stratified regression. Instead, subgroup-specific fits are used to assess whether the pooled model preserves coefficient direction, coefficient magnitude, and predictive performance across groups.

\section*{Approach}

The study combines a controlled simulation with one empirical illustration. The simulation varies the number of groups, total sample size, group-size imbalance, covariate distributions, and the degree of coefficient heterogeneity. All four scenarios use the same $140$-setting grid and $200$ replications per setting. Four settings with fewer than three observations in at least one subgroup are excluded from aggregate comparisons, leaving $136$ directly comparable settings per scenario.

The analysis compares pooled, group-adjusted, interaction, and separate subgroup regressions. It evaluates directional agreement, standardized coefficient movement, a covariate-context adjustment to that movement, the within-group loss from using the pooled rather than group-specific model, and covariate-distribution separation. The guiding distinction is between apparent disagreement created by finite samples or different predictor distributions and disagreement that remains when subgroup regression relationships genuinely differ.

Mathematically, if subgroup slopes differ, the pooled population target is a design-weighted compromise:
\[
\beta_{\mathrm{pool}}^{\dagger}
= \left(\sum_{g=1}^{G}X_g^\top X_g\right)^{-1}
  \sum_{g=1}^{G}X_g^\top X_g\beta_g.
\]
Even when all $\beta_g$ are equal, subgroup estimates vary because their precision depends on $X_g^\top X_g$. When slopes differ, the excess target-group risk from using the pooled rather than group-specific coefficient is proportional to
\[
(\beta_{\mathrm{pool}}^{\dagger}-\beta_g)^\top\Sigma_g(\beta_{\mathrm{pool}}^{\dagger}-\beta_g).
\]

\section*{Main simulation findings}

The unified simulation separates three interpretable patterns. Scenario 1 is the finite-sample baseline. Even with no true heterogeneity, splitting observations into many uneven groups increases subgroup coefficient dispersion. Scenario 2 changes predictor distributions while holding slopes fixed. It produces much larger covariate separation and substantial reduction after context normalization, but little additional prediction loss. Scenario 3 changes slopes while holding predictor distributions broadly comparable. It produces large residual coefficient movement and a much larger within-group loss from using the pooled model. Scenario 4 combines both mechanisms.

\begin{center}
\small
\begin{tabular}{lrrrrr}
\toprule
Scenario & Drift & Context drift & Prediction gap & Covariate shift & Shrinkage \\
\midrule
No heterogeneity & 3.64 & 3.11 & 0.004 & 0.057 & 5.4\% \\
Covariate shift only & 3.78 & 2.47 & 0.004 & 0.502 & 33.2\% \\
Coefficient heterogeneity only & 24.64 & 21.74 & 0.064 & 0.059 & 5.5\% \\
Combined heterogeneity & 25.41 & 16.12 & 0.065 & 0.508 & 33.5\% \\
\bottomrule
\end{tabular}
\end{center}

Relative to the no-heterogeneity baseline, coefficient heterogeneity increases median drift by $6.76$ times and the median prediction gap by $14.66$ times. The combined scenario has the largest raw drift and prediction gap. Its lower context-normalized drift than the coefficient-only case is itself informative: part of its raw movement is associated with covariate separation, but a large residual discrepancy remains.

On a reduced simulation grid, the conventional $\alpha=0.05$ interaction F-test rejected common slopes in $5.1\%$ of no-heterogeneity settings and $4.9\%$ of covariate-shift-only settings. Its average rejection rate was $85.2\%$ under coefficient heterogeneity and $84.5\%$ under combined heterogeneity. Thus, the standard test detects slope differences in this design, while the diagnostic framework adds information about coefficient magnitude, covariate separation, and the predictive cost of pooling.

\section*{Empirical illustration and contribution}

The empirical analysis uses substantively defined age groups in a health-survey regression. The pooled model masks a predictor whose coefficient changes direction in later age groups. The pattern remains after a sensitivity analysis excluding the smallest group, and group-specific models improve in-group predictive performance across the age strata. This is an illustration, not a causal claim or a universal biomedical conclusion.

Repeated five-fold cross-validation supports the in-sample comparison. Across 20 repetitions and 800 held-out group-fold evaluations, pooled RMSE exceeded age-specific RMSE by $0.911$ on average. The descriptive repeated-CV interval was $[0.813, 1.008]$, and the age-specific model had lower average held-out RMSE in every age group.

The contribution is therefore diagnostic and empirical. The paper provides a structured way to ask whether subgroup regression results justify a single pooled interpretation, and to distinguish three cases: finite-sample disagreement, disagreement associated with covariate context, and persistent relationship heterogeneity. The intended use is as a practical complement to interaction models and stratified regression, not as a substitute for either.

\section*{Current limits and request}

The work is limited to low-dimensional linear models, a restricted set of simulated mechanisms, and one empirical illustration. The context adjustment is descriptive rather than a formal causal decomposition. The $200$ replications per setting provide substantially more stable Monte Carlo summaries than the original exploratory run. The interaction benchmark uses a reduced grid, and the empirical validation uses repeated held-out comparisons from one dataset. Additional datasets and alternative subgroup definitions remain important extensions.

The manuscript is now internally consistent, fully compiled, and supported by a unified four-scenario simulation design. I am seeking feedback on whether the framing and evidence are appropriate for posting as a cautious preprint, and whether the contribution is sufficiently clear to support an endorsement.

\end{document}